%-----------------------------------------------------------------------------------------------------------------------------------------------%
%	The MIT License (MIT)
%
%	Copyright (c) 2015 Jan Küster
%
%	Permission is hereby granted, free of charge, to any person obtaining a copy
%	of this software and associated documentation files (the "Software"), to deal
%	in the Software without restriction, including without limitation the rights
%	to use, copy, modify, merge, publish, distribute, sublicense, and/or sell
%	copies of the Software, and to permit persons to whom the Software is
%	furnished to do so, subject to the following conditions:
%	
%	THE SOFTWARE IS PROVIDED "AS IS", WITHOUT WARRANTY OF ANY KIND, EXPRESS OR
%	IMPLIED, INCLUDING BUT NOT LIMITED TO THE WARRANTIES OF MERCHANTABILITY,
%	FITNESS FOR A PARTICULAR PURPOSE AND NONINFRINGEMENT. IN NO EVENT SHALL THE
%	AUTHORS OR COPYRIGHT HOLDERS BE LIABLE FOR ANY CLAIM, DAMAGES OR OTHER
%	LIABILITY, WHETHER IN AN ACTION OF CONTRACT, TORT OR OTHERWISE, ARISING FROM,
%	OUT OF OR IN CONNECTION WITH THE SOFTWARE OR THE USE OR OTHER DEALINGS IN
%	THE SOFTWARE.
%	
%
%-----------------------------------------------------------------------------------------------------------------------------------------------%


%============================================================================%
%
%	DOCUMENT DEFINITION
%
%============================================================================%

%we use article class because we want to fully customize the page and don't use a cv template
\documentclass[10pt,A4]{article}	


%----------------------------------------------------------------------------------------
%	ENCODING
%----------------------------------------------------------------------------------------

%we use utf8 since we want to build from any machine
\usepackage[utf8]{inputenc}		

%----------------------------------------------------------------------------------------
%	LOGIC
%----------------------------------------------------------------------------------------

% provides \isempty test
\usepackage{xifthen}

%----------------------------------------------------------------------------------------
%	FONT
%----------------------------------------------------------------------------------------

% some tex-live fonts - choose your own

%\usepackage[defaultsans]{droidsans}
%\usepackage[default]{comfortaa}
%\usepackage{cmbright}
%\usepackage[default]{raleway}
%\usepackage{fetamont}
%\usepackage[default]{gillius}
%\usepackage[light,math]{iwona}
%\usepackage[thin]{roboto} 

% set font default
\renewcommand*\familydefault{\sfdefault} 	
\usepackage[T1]{fontenc}

% more font size definitions
\usepackage{moresize}		


%----------------------------------------------------------------------------------------
%	PAGE LAYOUT  DEFINITIONS
%----------------------------------------------------------------------------------------

%debug page outer frames
%\usepackage{showframe}			


%define page styles using geometry
\usepackage[a4paper]{geometry}		

% for example, change the margins to 2 inches all round
\geometry{top=1.25cm, bottom=-.6cm, left=1.5cm, right=1.5cm} 	

%use customized header
\usepackage{fancyhdr}				
\pagestyle{fancy}

%less space between header and content
\setlength{\headheight}{-5pt}		


%customize entries left, center and right
%\lhead{}
%\chead{}
%\rhead{}


%indentation is zero
\setlength{\parindent}{0mm}

%----------------------------------------------------------------------------------------
%	TABLE /ARRAY DEFINITIONS
%---------------------------------------------------------------------------------------- 

%for layouting tables
\usepackage{multicol}			
\usepackage{multirow}

%extended aligning of tabular cells
\usepackage{array}

\newcolumntype{x}[1]{%
>{\raggedleft\hspace{0pt}}p{#1}}%


%----------------------------------------------------------------------------------------
%	GRAPHICS DEFINITIONS
%---------------------------------------------------------------------------------------- 

%for header image
\usepackage{graphicx}

%for floating figures
\usepackage{wrapfig}
\usepackage{float}
%\floatstyle{boxed} 
%\restylefloat{figure}

%for drawing graphics		
\usepackage{tikz}				
\usetikzlibrary{shapes, backgrounds,mindmap, trees}


%----------------------------------------------------------------------------------------
%	Color DEFINITIONS
%---------------------------------------------------------------------------------------- 

\usepackage{color}

%accent color
\definecolor{sectcol}{RGB}{0,0,180}

%dark background color
\definecolor{bgcol}{RGB}{110,110,110}

%light background / accent color
\definecolor{softcol}{RGB}{225,225,225}


%----------------------------------------------------------------------------------------
%	Other Packages
%---------------------------------------------------------------------------------------- 

% for \href
\definecolor{linkcolor}{RGB}{0,50,150}
\usepackage[colorlinks=true,urlcolor=linkcolor,linkcolor=linkcolor]{hyperref}

%============================================================================%
%
%
%	DEFINITIONS
%
%
%============================================================================%

%----------------------------------------------------------------------------------------
% 	HEADER
%----------------------------------------------------------------------------------------

% remove top header line
\renewcommand{\headrulewidth}{0pt} 

%remove bottom header line
\renewcommand{\footrulewidth}{0pt}	  	

%remove pagenum
\renewcommand{\thepage}{}	

%remove section num		
\renewcommand{\thesection}{}			

%----------------------------------------------------------------------------------------
% 	ARROW GRAPHICS in Tikz
%----------------------------------------------------------------------------------------

% a six pointed arrow pointing to the left
\newcommand{\tzlarrow}{(0,0) -- (0.2,0) -- (0.3,0.2) -- (0.2,0.4) -- (0,0.4) -- (0.1,0.2) -- cycle;}	

% include the left arrow into a tikz picture
% param1: fill color
%
\newcommand{\larrow}[1]
{\begin{tikzpicture}[scale=0.58]
	 \filldraw[fill=#1!100,draw=#1!100!black]  \tzlarrow
 \end{tikzpicture}
}

% a six pointed arrow pointing to the right
\newcommand{\tzrarrow}{ (0,0.2) -- (0.1,0) -- (0.3,0) -- (0.2,0.2) -- (0.3,0.4) -- (0.1,0.4) -- cycle;}

% include the right arrow into a tikz picture
% param1: fill color
%
\newcommand{\rarrow}
{\begin{tikzpicture}[scale=0.7]
	\filldraw[fill=sectcol!100,draw=sectcol!100!black] \tzrarrow
 \end{tikzpicture}
}



%----------------------------------------------------------------------------------------
%	custom sections
%----------------------------------------------------------------------------------------

% create a coloured box with arrow and title as cv section headline
% param 1: section title
%
\newcommand{\cvsection}[1]
{
	\begin{center}
		\large\textcolor{sectcol}{\textbf{#1}}
	\end{center}
}

%create a coloured arrow with title as cv meta section section
% param 1: meta section title
%
\newcommand{\metasection}[2]
{
%\begin{tabular*}{1\textwidth}{r r}
\footnotesize{#2} \hspace*{\fill} \footnotesize{#1}\\[1pt]
%\end{tabular*}
}

%----------------------------------------------------------------------------------------
%	 CV EVENT
%----------------------------------------------------------------------------------------

% creates a stretched box as cv entry headline followed by some paragraphs about 
% the work you did
% param 1:	event time i.e. 2014 or 2011-2014 etc.
% param 2:	event name (what did you do?)
% param 3:	institution (where did you work / study)
% param 4:	list of paragraphs
%
\newcommand{\cvevent}[4]
{

\begin{tabular*}{1\textwidth}{p{13.6cm}  x{3.9cm}}
	\textbf{#2} - \textcolor{bgcol}{#3} &   \vspace{2.5pt}\textcolor{sectcol}{#1}
\end{tabular*}

\vspace{-8pt}
\textcolor{softcol}{\hrule}
\vspace{6pt}

	\foreach \desc in {#4}{
		$\cdot$ \desc\\[3pt]
	}
	
\vspace{3pt}
}

% creates a stretched box as 
\newcommand{\cveventmeta}[2]
{
	\mbox{\mystrut \hspace{87pt}\textit{#1}}\\
	#2
}

%----------------------------------------------------------------------------------------
% CUSTOM STRUT FOR EMPTY BOXES
%----------------------------------------- -----------------------------------------------
\newcommand{\mystrut}{\rule[-.3\baselineskip]{0pt}{\baselineskip}}

%----------------------------------------------------------------------------------------
% CUSTOM LOREM IPSUM
%----------------------------------------------------------------------------------------
\newcommand{\lorem}
{Lorem ipsum dolor sit amet, consectetur adipiscing elit. Donec a diam lectus.}



%============================================================================%
%
%
%
%	DOCUMENT CONTENT
%
%
%
%============================================================================%
\begin{document}

%use our custom fancy header definitions
\pagestyle{fancy}

%---------------------------------------------------------------------------------------
%	TITLE HEADLINE
%----------------------------------------------------------------------------------------

\vspace{-8pt}
\begin{center}
    {\HUGE \textsc{Boudrouss Réda} \textcolor{sectcol}{\rule[-1mm]{1mm}{0.9cm}} \textsc{CV}}\\[2pt]
    {\small Candidat en thèse \textbf{CIFRE} -- Analyse statique \& interprétation abstraite -- Paris}
\end{center}

\vspace{6pt}

%----------------------------------------
%   MÉTADONNÉES
%----------------------------------------
\metasection{\href{https://github.com/rboudrouss}{github.com/rboudrouss}}{}
\metasection{\href{https://www.linkedin.com/in/rboudrouss}{linkedin.com/in/rboudrouss}}{}
\metasection{\href{mailto:contact@rboud.com}{contact@rboud.com}}{}
\metasection{\href{https://rboud.com}{rboud.com}}{}
\metasection{+33 (0)7 55 78 85 16}{\textbf{Langues :} Français (Natif), Anglais (Professionnel), Arabe (Conversationnel)}

\vspace{-15pt}
\textcolor{softcol}{\hrule}

\normalsize

%----------------------------------------
%   RÉSUMÉ
%----------------------------------------
\cvsection{Résumé}
En M2 Sciences et Technologie du Logiciel à Sorbonne Université, je recherche une \textbf{entreprise partenaire pour une thèse CIFRE} en \textbf{analyse statique}, encadrée par \href{https://mine.perso.lip6.fr/index.html.fr}{\textbf{Antoine Miné}} (LIP6) : \textbf{interprétation abstraite} appliquée à \textbf{WebAssembly}, \textbf{Rust} et aux langages dynamiques comme \textbf{JavaScript}. J'ai porté l'analyseur \textbf{MOPSA} (OCaml/C/C++) vers WebAssembly et écrit un analyseur par interprétation abstraite en Rust. Mon alternance chez Hexaglobe m'a donné une pratique concrète de la \textbf{vérification logicielle en contexte industriel}.

%----------------------------------------
%   FORMATION
%----------------------------------------
\cvsection{Formation}

\cvevent{2024 -- 2026}{Master Science et Technologie du Logiciel}{Sorbonne Université}{
    {Cours suivis : \textbf{interprétation abstraite} et systèmes de types (cours d'Antoine Miné), inférence de propriétés numériques, \textbf{vérification formelle}, typage statique, conception et implémentation de \textbf{langages de programmation} et de leurs environnements d'exécution (interpréteurs, compilateurs), \textbf{concurrence} et programmation répartie.}
}

\cvevent{2020 -- 2024}{Licence Informatique}{Sorbonne Université}{
    {Algorithmique, architecture des ordinateurs, systèmes d'exploitation, réseaux, bases de données, programmation fonctionnelle.}
}

%----------------------------------------
%   PROJETS
%----------------------------------------
\vspace{-22pt}
\cvsection{Projets}

\cvevent{2026}{Mopsa-wasm}{\href{https://github.com/rboudrouss/mopsa-emcc}{github.com/rboudrouss/mopsa-emcc}}{
	{Portage de \textbf{MOPSA}, analyseur statique par \textbf{interprétation abstraite} pour C et Python écrit en \textbf{OCaml}/C/C++ au LIP6, vers \textbf{WebAssembly} : compilation Emscripten du runtime OCaml et de toute la pile native (frontend Clang, \textbf{Apron}, GMP, MPFR) en un module Wasm statique unique.},
	{Travail des deux côtés de la frontière \textbf{runtime OCaml / FFI C} : table statique d'environ 1 400 primitives C (pas de \texttt{dlopen} en Wasm), correction d'un bug de dépassement non signé dans la macro \texttt{Tag\_val}. Traitement du risque de \textbf{correction} lié au mode d'arrondi figé de Wasm en recompilant les domaines numériques d'Apron sur des rationnels exacts GMP. Publication d'un \href{https://mopsawasm.rboud.com}{playground} côté client et d'un \href{https://rboud.com/blog/running-mopsa-in-the-browser/}{article technique}.}
}

\cvevent{2026}{Reactant analyzer}{\href{https://github.com/rboudrouss/reactant-analyzer}{github.com/rboudrouss/reactant-analyzer}}{
	{Conception et implémentation d'un analyseur statique par \textbf{interprétation abstraite} en \textbf{Rust} (environ 30k lignes, 500+ tests) détectant des bugs de hooks \textbf{React} hors de portée des linters syntaxiques : boucles de rendu infinies, état dérivé, dépendances instables.},
	{Évaluation des composants dans des domaines abstraits (intervalles, ensembles de constantes, étiquettes de stabilité référentielle) sur une IR dédiée à base de CFG  : une boucle de rendu infinie est signalée lorsque le point fixe sur l'état \textbf{élargit}. Sémantique concrète fondée sur l'article React-tRace (OOPSLA 2025).}
}

%----------------------------------------
%   EXPÉRIENCES
%----------------------------------------
\vspace{-22pt}
\cvsection{Expériences}

\cvevent{2025 -- 2026}{Quality Analyst Engineer (alternance)}{Hexaglobe (streaming vidéo)}{
	{Conception de \textbf{stratégies de test} complètes (unitaires, intégration, E2E) pour des produits legacy sans aucune couverture, jusqu'à environ \textbf{3 000 cas de test} sur 4 stacks (TypeScript, Python, PHP, JS).},
	{Développement d'\textbf{oracles de qualité objectifs} (VMAF, détection de trames noires et de gels d'image, analyse audio) pour valider un encodeur vidéo live.},
	{Industrialisation d'une production de tests par \textbf{agents IA} conditionnée à des signaux de vérification déterministes (orchestration multi-agents en sandbox Docker, garde-fous versionnés, tests verts obligatoires, relecture humaine) : environ 720 tests d'intégration en 3 jours.},
	{Pipelines \textbf{GitLab CI/CD} et composants CI mutualisés utilisés par 10+ dépôts, avec intégration d'outils d'\textbf{analyse statique} dans la chaîne.}
}

\cvevent{2023 -- 2026}{Responsable technique, puis Président (2024--25)}{Association Play Sorbonne Université}{
	{Conception et exploitation d'une \textbf{infrastructure virtualisée} (Docker, CI/CD, réseau) et développement d'une web app mobile \textbf{local-first}. Comme président, encadrement de \textbf{100+ bénévoles} pour un événement accueillant 10 000+ visiteurs.}
}

%--------------------------------------------------------------------------------------------------
%	ARTIFICIAL FOOTER (fancy footer cannot exceed linewidth)
%--------------------------------------------------------------------------------------------------

\null
\vspace*{\fill}
\hspace{-0.25\linewidth}\colorbox{white}{\makebox[1.5\linewidth][c]{\mystrut  \textnormal{\textcolor{sectcol}{\href{https://www.linkedin.com/in/rboudrouss}{linkedin.com/in/rboudrouss}} $\cdot$ \textcolor{sectcol}{\href{https://github.com/rboudrouss}{github.com/rboudrouss}}}}}

\end{document}
